% !PRINTED 446
\weberpage{446}
Nach \S.\,115, (14) und (8) ist aber

\[
\alpha^2 + \beta^2\varepsilon^3 + \alpha\,\beta\,\varepsilon^2 = \alpha^2 + \beta\,(\alpha\,\varepsilon^2 + \beta\,\varepsilon^3) = \alpha^2 - \beta\,\gamma = -\,\alpha,
\]
und wenn man die drei Co\"efficienten des Ausdrucks (8) in dieser Weise umformt, so ergiebt sich

\[
-\,\alpha\,\gamma\,x_1 - \beta\,\gamma\,\varepsilon^3 x_2 - \alpha\,\beta\,\varepsilon^2 x_3,
\]
d. h. $A$ \"andert durch Anwendung der Substitution $\omega$ sein Vorzeichen, und die Hauptaxe $A$ bleibt also durch die Substitution $\omega$ unge\"andert.

Da $A$ ausserdem durch die Substitution $\Theta$ unge\"andert bleibt, weil zwei Pole von $\Theta$ auf $A$ liegen, so haben wir die ganze Gruppe 8\textsuperscript{ten} Grades, durch die $A$ unge\"andert bleibt, die wir die \gesperrt{Gruppe von} $A$ nennen und mit $G_a$ bezeichnen, in der Form

\weq{9}{\omega^u\,\Theta^v.}
Unter diesen acht Substitutionen sind nur zwei, n\"amlich die identische und $\Theta^2$, die alle Punkte der Linie $A$ unge\"andert lassen. Bei den anderen bleiben nur je zwei Punkte in Ruhe. Denn berechnet man auf dem hier eingeschlagenen Wege aus \S.\,115, (5), (16) die Hauptaxen von $\omega$, $\omega\Theta$, $\omega\Theta^2$, $\omega\Theta^3$, so erh\"alt man, in der Bezeichnung (7),

\[
A_{2,1},\ A_{3,0},\ A_{3,3},\ A_{2,0},
\]
die alle von der Hauptaxe $A_{1,0}$ von $\Theta^2$ verschieden sind.

Die Substitution $\Theta^2$ hat ausser der Hauptaxe $A$ noch einen isolirten Pol, den wir erhalten, wenn wir die Wurzel $\lambda = 1$ der cubischen Gleichung (4) w\"ahlen. Dann ergeben sich f\"ur die Coordinaten dieses Poles die Gleichungen

\[
\begin{array}{l}
(\beta-1)\,x_1 + \gamma\,\varepsilon^3 x_2 + \alpha\,\varepsilon^2 x_3 = 0,\\
\gamma\,\varepsilon^{-3} x_2 + (\alpha-1)\,x_2 + \beta\,\varepsilon^{-1} x_3 = 0,\\ % [sic] printed x_2 twice, expected x_1 in first term
\alpha\,\varepsilon^{-2} x_1 + \beta\,\varepsilon\, x_2 + (\gamma-1)\,x_3 = 0,
\end{array}
\]
und daraus erh\"alt man

\weq{10}{x_1 : x_2 : x_3 = \alpha\,\gamma\,\varepsilon^4 : \beta\,\gamma\,\varepsilon : \alpha\,\beta\,\varepsilon^2.}
Solcher Punkte giebt es 21. Den Punkt (10) bezeichnen wir f\"ur den Augenblick mit $a$. Der Punkt $a$ bleibt nun nicht nur durch $\Theta$, sondern auch, wie eine Rechnung auf Grund der Formeln \S.\,115, (17), (14) zeigt, durch $\omega$, und folglich durch die ganze Gruppe $G_a$ unge\"andert, und ist also ein mindestens achtz\"ahliger Pol. Da er aber durch 21 Substitutionen in 21 verschiedene Lagen gebracht werden kann, so ist er auch nicht \gesperrt{mehr} als
%CONT>>
